\begin{claim}{C.36}{Proved}{Compatible contractions on two-dimensional proper closures}
Let $f_t$ be a based ambient-null plane loop, let $G$ be
its generated group, and let $H$ be any group of plane
homeomorphisms commuting with every $f_t$.
For an open $G$-orbit $O$, let $Q_O$ be the closure of
the restricted setwise $H$-stabilizer, as in C.30.
For every such $O$ with $\dim Q_O=2$, there is a based
null-homotopy on the entire closure $\overline O$
through homeomorphisms of $\overline O$ that commute
with its operator stabilizer.
These closure contractions agree on every intersection
$\overline O_1\cap\overline O_2$ and are carried to one
another by $H$. They therefore give a jointly continuous
family, with jointly continuous inverse, on every
finite union of these closures.
No continuity across an infinite accumulating family
of closures is asserted.
\end{claim}
\begin{proof}
\emph{The actual quotient action on the closure.}
C.30 makes $Q_O$ a closed locally compact free proper
group of homeomorphisms of $O$, and C.33 makes it Lie.
Its dimension-two argument shows that $Q_O$ is
connected and simply transitive on $O$, and that
\[
 N_O=C_{\operatorname{Homeo}(O)}(Q_O)
                    \cong Q_O^{\mathrm{op}}             \tag{C.36.1}
\]
is also a connected, simply transitive two-dimensional
Lie group. Evaluation $N_O\to O$ is a homeomorphism.
The image of $G$ in $N_O$ is all of $N_O$: for any
$n\in N_O$, transitivity of $G$ supplies $g\in G$
with $gx=nx$ at one point $x\in O$, and the
injectivity of evaluation on $N_O$ then gives
$g|_O=n$.

Put $L=\overline G$ in $\operatorname{Homeo}^+(\R^2)$.
This is a Polish group. It preserves every open
$G$-orbit $O$. Indeed, if $g_j\to l$ with $g_j\in G$,
the identities $g_j(\partial O)=\partial O$ and
the convergent inverse family give
$l(\partial O)=\partial O$. For $x\in O$,
$lx\in\overline O$ by convergence, and
$lx\notin\partial O$ because $l$ preserves that
boundary and is injective. Thus $lx\in O$;
applying the same argument to $l^{-1}$ proves
$l(O)=O$.
The elements of $L$ still commute with $H$.
Restriction consequently gives a continuous
surjective homomorphism
\[
                   r_O:L\longrightarrow N_O.            \tag{C.36.2}
\]
It is open by the open mapping theorem for Polish
groups. Equivalently, apply the open orbit-map
principle from
\cite[Theorem~3.5]{C_GomezPinedoUzcategui2017}
to the transitive action of $L$ on the Baire
space $N_O$ through (C.36.2).

The kernel of (C.36.2) fixes $O$ pointwise, and
hence fixes $\overline O$ pointwise. Therefore
there is a well-defined action
\[
 N_O\times\overline O\longrightarrow\overline O,
       \qquad(n,z)\longmapsto lz
                 \quad\text{for any }l\text{ with }r_O(l)=n.
                                                        \tag{C.36.3}
\]
This action is jointly continuous: $r_O\times\id$
is an open quotient map, and the continuous
$L$-action factors through it. No continuous
section of (C.36.2) is required.
In fact each $n$ has a preimage in $G$, as proved
above. Hence (C.36.3) preserves every original
$G$-orbit contained in $\overline O$, and its
image on such an orbit is exactly the image of
$G$. This conclusion does not require a claim
that $L$ preserves every other $G$-orbit in the
plane.

\emph{The two-dimensional Lie contraction.}
We record the Lie groups and the maps being used.
A two-dimensional real Lie algebra is either
abelian or has a basis $X,Y$ with $[X,Y]=Y$:
in the nonabelian case its one-dimensional
derived algebra is an ideal, and scaling a
complementary basis vector gives that bracket.
The simply connected groups are respectively
$\R^2$ and the positive affine group of the
line. In coordinates $(a,b)$, the latter has
product
\[
           (a,b)(a',b')=(a+a',\,b+e^a b').
                                                        \tag{C.36.4}
\]
Its exponential map is
\[
 \exp(u,v)=\left(u,v\,\frac{e^u-1}{u}\right),
 \quad\text{with }\frac{e^u-1}{u}=1\text{ at }u=0.
                                                        \tag{C.36.5}
\]
The scalar factor is positive for every real
$u$, so this is a global diffeomorphism.
For $\R^2$ the exponential is the identity.
The affine group has trivial center, so its
connected groups with this Lie algebra have
no further covering quotients. In the abelian
case the connected group is a quotient of
$\R^2$ by a discrete central subgroup.
The compact torus is excluded here because
$N_O\cong O$ is noncompact.
Thus in every case the universal covering
group $\widetilde N_O$ has a globally defined
inverse $\log$ to its exponential.

Let $n_O(t)=f_t|_O\in N_O$.
By the restriction assertion of C.32 this loop
is null-homotopic in $\operatorname{Homeo}(O)$.
Evaluation on $O$ and the homeomorphism
$N_O\to O$ show that it is null-homotopic
in $N_O$. Its unique based lift
$\widetilde n_O(t)$ to $\widetilde N_O$
is therefore a closed loop.
With $p_O:\widetilde N_O\to N_O$ the covering
homomorphism, define
\[
 R^O_{s,t}
   =p_O\!\left(\exp\bigl((1-s)\log\widetilde n_O(t)\bigr)\right)
       \quad\text{acting on }\overline O
       \text{ by (C.36.3)}.                             \tag{C.36.6}
\]
This is a continuous family of homeomorphisms
with continuous inverse, starts with $f_t$ on
$\overline O$, ends with the identity, and
keeps both based edges fixed.
On $O$ it lies in the exact centralizer of
$Q_O$. Its commutation with the original
setwise operator stabilizer extends to
$\overline O$ by continuity and density.

\emph{The common boundary image and its topology.}
Take two of these closures and a $G$-orbit
$B$ in their intersection. It is a point,
an intrinsic line, or a circle: an open orbit
cannot lie in the boundary of another open
orbit. The point case is fixed throughout.
For a nontrivial $B$, each $N_i=N_{O_i}$
acts transitively on $B$ by (C.36.3), and
both action images are exactly $G|_B$.
We verify that this common image has a single
Lie topology, rather than identify abstract
quotients without their topology.

For $p\in B$, its stabilizer in $N_i$ is
closed. The homogeneous space
$N_i/(N_i)_p$ maps continuously and
bijectively to $B$. Its dimension is one:
dimension zero would make the connected
orbit a point, and dimension two would make
$B$ open in the plane by invariance of
domain. Its topology agrees with the
intrinsic topology of $B$.
Indeed, a compact interval in this
homogeneous one-manifold maps to an
embedded compact arc. In an intrinsic
line orbit that arc is the unique
connecting arc supplied by C.28; on a
circle the same conclusion holds locally.
The interiors of these arcs give
intrinsic interval neighborhoods, so the
bijective orbit map is a local, and
therefore global, intrinsic homeomorphism.
Thus the action is continuous also in
the intrinsic topology.

Divide $N_i$ by its closed kernel on $B$.
The resulting connected Lie group $M_i$
acts faithfully and transitively on that
one-manifold, and has dimension at most
two. In dimension one it is abelian, so
a point stabilizer is normal and is
therefore the kernel; faithfulness makes
the stabilizer trivial. Evaluation then
identifies its action with the translation
group of $\R$ or the rotation group of
$\R/\Z$.
An abelian group of dimension two would
give the same freeness conclusion and
an orbit of dimension two, so cannot
occur faithfully here.
In the remaining case $M_i$ is the
positive affine group (C.36.4).
The identity component of a point
stabilizer is a one-dimensional subgroup.
It cannot be the translation subgroup,
which is normal and would lie in the
action kernel. Every other such subgroup
is conjugate to the dilation subgroup.
Its normalizer is that dilation subgroup
itself, as is seen by conjugating its
unique fixed point on the line.
The full point stabilizer is consequently
the dilation subgroup, and the homogeneous
action is the usual positive affine
action on $\R$.

These descriptions prove topological
embedding of each actual image in
$\operatorname{Homeo}(B)$ with its
intrinsic compact-open topology.
For translations or rotations one
evaluation recovers the parameter.
For the affine group, the images of
two distinct ordered points recover
the positive slope and the intercept
continuously. They also prove that
the image groups are closed in that
homeomorphism group.
Since both images are the same set
$G|_B$, we have obtained a common
topological Lie group $M_B$, and
continuous quotient homomorphisms
\[
                       \rho_i:N_i\longrightarrow M_B.
                                                        \tag{C.36.7}
\]
The universal cover of $M_B$ is
$\R$ or the positive affine group,
and its exponential is again a
global diffeomorphism.

\emph{Equality on overlaps.}
The homomorphisms (C.36.7) lift
uniquely to based homomorphisms
\[
               \widetilde\rho_i:
                    \widetilde N_i\longrightarrow
                    \widetilde M_B.
                                                        \tag{C.36.8}
\]
The two paths
$\widetilde\rho_i(\widetilde n_i(t))$
are based lifts of the identical
loop $f_t|_B$, so covering uniqueness
makes them equal to one path
$\widetilde b(t)$.
Naturality of the Lie exponential
and its global injectivity give
\[
 \begin{split}
 d\widetilde\rho_i(\log\widetilde n_i(t))
      &=\log\widetilde b(t),\\
 \widetilde\rho_i
   \!\left(\exp\bigl((1-s)\log\widetilde n_i(t)\bigr)\right)
      &=\exp\bigl((1-s)\log\widetilde b(t)\bigr).
 \end{split}                                             \tag{C.36.9}
\]
After projection to $M_B$, this says
precisely that $R^{O_1}_{s,t}$ and
$R^{O_2}_{s,t}$ agree on $B$.
The intersection of the two closures
is a union of such $G$-orbits, so
they agree everywhere on the
intersection. This includes circle
orbits with their actual based
universal-cover lifts; no choice
of an unrecorded integer rotation
is made.

For $h\in H$, conjugation carries
$Q_O$ onto $Q_{hO}$ and $N_O$
onto $N_{hO}$, as follows from
the exact restriction of the
setwise stabilizers. This is a
topological Lie-group isomorphism.
It carries $n_O(t)$ to
$n_{hO}(t)$, since $h$ commutes
with $f_t$. Its based lift and
exponential naturality show
\[
                    R^{hO}_{s,t}
                         =hR^O_{s,t}h^{-1}.
                                                        \tag{C.36.10}
\]
Thus the compatible families have
the asserted operator equivariance.

Finally, on a finite union of
these closures, their continuous
maps agree on every overlap.
The finite closed-set pasting
lemma gives joint continuity.
Each intersection is preserved,
because its $G$-orbits are
preserved; hence the inverse
maps also agree there.
Applying the same pasting
argument to the inverse families
gives the stated family of
homeomorphisms on the finite
union. This proves the claim.
\end{proof}

\paragraph{Exact scope of the closure construction.}
The quotients (C.36.2)--(C.36.3)
provide continuity on each
individual closure, and
(C.36.7)--(C.36.9) enforce the
actual common boundary restriction
relations. In particular they
respect the three dependent edge
characters in C.35.
The present proof uses finite
closed-set pasting. C.37
subsequently proves a uniform
three-word estimate and extends
these contractions continuously,
with their inverses, to the
closed union $X_2$ of all such
open orbits, including infinite
accumulation. Extension across
$X_2$ to its exterior, including
frontier contractions for
$\dim Q_O=0,1$ and compatible
choices elsewhere on $Y$,
remains required for a global
conclusion in C.12.
